\section{Einleitung}

Der im Rahmen des Projektstudiums im Fernstudium Informatik zu erstellende Geometry Wars-Klon ist in~\cite{SH25} beschrieben.
Wir gehen im Folgenden auf einige Zieldaten ein, um sie im späteren Verlauf referenzieren und vergleichen zu können.\\

\noindent
Als zu realisierende Features wurden vereinbart:

\begin{itemize}
    \itemsep0.5em
    \item Spielbares Level (2D-Grid)
    \item Time-Attack Modus mit einer Dauer von 3 Minuten
    \item Highscore-System
    \item Controller-Steuerung
    \item Stabile Framerate von mindestens 60 FPS bei gleichzeitiger Darstellung von über hundert Gegnern und Objekten
    \item Drei Gegnertypen mit einfacher KI (``Spawn and Chase``)
    \item Aufwertung des Game Feels durch grafische Effekte
    \item Kompilierbar und lauffähig unter Windows 11
\end{itemize}

\noindent
Mit Veröffentlichung des o.a. Dokumentes war bereits Meilenstein 1 erreicht, welcher den Fokus auf das außenliegende Gerüst der Anwendung gelegt hat. \\
In~\ref{tab:zeitplan} ist eine Übersicht der nachfolgenden Meilensteine 1 - 6 dargestellt, deren Fertigstellungstermine mit einer jeweiligen Verzögerung von 2 - 4 Wochen erreicht worden sind.
Hier hat sich gezeigt, dass mit der zunehmenden Komplexität der umzusetzenden Features wie Input-System, Spawn-Systeme, Kollisionsabfrage und UI-Rendering auch immer wieder versteckte Anforderungen Anpassungen bereits existierender Systeme erfordert habe.
So hat bspw, das Input-System auch die Berücksichtigung von Stick-Drift erfordert; Spawning von Gegnerobjekten hat ein Pool von wiederverwendbaren Objekten (\textit{prefabricated objects}, i.F. \textit{prefab}) erfordert; zusätzlich zu der Erkennung von Kollisionen haben sich komplexe Fragestellungen hinsichtlich der Verarbeitung der Kollisionen mit dem beteiligten Systemen ergeben; UI-Rendering hat eine Erweiterung der Rendering Engine um die Unterstützung von Texturen und UI-Elementen erfordert, sowie die Implementierung von hierarchischer endlicher Automaten, um das Spiel in verschiedene Zustände (z.B. ``Main Menu``, ``In Game``, ``Game Over``) sauber zu kapseln.\\



\setlength{\tabcolsep}{8pt}
\begin{table}[t]
    \centering
    {\renewcommand{\arraystretch}{1.2}%
        \begin{tabularx}{\textwidth}{@{} l l Y @{}}
        \toprule
        \textbf{Meilenstein} & \textbf{Datum} & \textbf{Inhalt} \\
        \midrule
        \texttt{milestone\_1} & 20.10.2025 &
        Bereitstellung der Applikationsschicht inkl. Implementierung des Event-Systems,
        des Input-Managers sowie Anbindung an das Low-Level-API-Subsystem. \\
        \texttt{milestone\_2} & 17.11.2025 &
        Bereitstellung der Rendering-Engine; erster Entwurf des Spielfelds samt Darstellung des Spielerschiffs. \\
        \texttt{milestone\_3} & 22.12.2025 &
        Umsetzung der Physik und Spielereingabe zur Steuerung des Schiffs; Feuermechaniken. \\
        \texttt{milestone\_4} & 19.01.2026 &
        Implementierung der wichtigsten Spielregeln und -mechaniken; spielfähiger Prototyp. \\
        \texttt{milestone\_5} & 09.02.2026 &
        Bereitstellung des Prototyps; Feinschliff. \\
        \texttt{milestone\_6} & 16.03.2026 &
        Abgabe der Dokumentation und Vorstellung des Projekts. \\
        \bottomrule
        \end{tabularx}}
    \caption{Geplante Meilensteine zur Umsetzung des \textit{Geometry Wars}-Klon.
    Zum Zeitpunkt der Veröffentlichung von~\cite{SH25} waren die Meilensteine 1-2 bereits abgeschlossen, Meilensteine 3-6 wurden mit leichter Verzögerung entsprechend der Anforderungen erreicht.}
    \label{tab:zeitplan}
\end{table}




